سیستم \lr{Transform Coding} داده‌های تصویر را به حوزهٔ تبدیل می‌برد تا انرژی در تعداد کمی ضریب متمرکز شده و فشرده‌سازی کارآمد شود. با ارجاع به اسلاید ۷۵ از \lr{Chapter} 11 درس که مربوط به \lr{Image Compression} است، مسیر کدگذار (سطر بالا) و مسیر کدگشا (سطر پایین) را گام‌به‌گام شرح می‌دهیم:

\subsection*{کدگذار (\lr{Encoder})}

با توجه به اسلاید ۷۵ داریم:

\lr{Construct $n \times n$ subimages}:
تصویر ورودی $N \times N$ به بلوک‌های کوچک و بدون هم‌پوشانی $n \times n$ افراز می‌شود تا هر بلوک مستقل تبدیل گردد؛ این کار همبستگی محلی پیکسل‌ها را بهره‌برداری کرده و هزینهٔ محاسباتی تبدیل را کاهش می‌دهد.

\lr{Forward transform}:
هر زیرتصویر از یک تبدیل برگشت‌پذیر (مانند \lr{DCT}) عبور می‌کند که پیکسل‌ها را ناهمبسته کرده و بخش عمدهٔ انرژی سیگنال را در چند ضریب کم‌فرکانس متمرکز می‌سازد.

\lr{Quantizer}:
ضرایب تبدیل کوانتیزه می‌شوند؛ این تنها گام اتلافی و برگشت‌ناپذیر زنجیره است. ضرایب کم‌اهمیت بصری (پرفرکانس و کم‌انرژی) درشت کوانتیزه یا صفر می‌شوند و بیشترین نرخ فشرده‌سازی از همین‌جا حاصل می‌گردد.

\lr{Symbol encoder}:
ضرایب کوانتیزه‌شده با کدگذاری آنتروپی (مانند \lr{run-length} و کد طول‌متغیر) کد می‌شوند تا افزونگی کدگذاری حذف شده و تصویر فشرده تولید گردد.

\subsection*{کدگشا (\lr{Decoder})}

\lr{Symbol decoder}:
کدگذاری آنتروپی معکوس شده و ضرایب کوانتیزه‌شده بازیابی می‌گردند.

\lr{Inverse transform}:
معکوس تبدیل مستقیم بر هر بلوک اعمال می‌شود و آن را به حوزهٔ مکانی (پیکسل) بازمی‌گرداند.

\lr{Merge $n \times n$ subimages}:
بلوک‌های بازسازی‌شده دوباره کنار هم چیده می‌شوند تا تصویر کامل از فشردگی خارج‌شده ساخته شود.

\subsection*{حالت موجک (\lr{WT})}

حال معماری اسلاید ۷۹ را بررسی ‌می‌نماییم.

در سیستم کدگذاری مبتنی بر تبدیل موجک، گام‌های ساخت و ادغام زیرتصاویر حذف می‌شوند؛ زیرا تبدیل موجک یک‌جا روی کل تصویر اعمال می‌گردد. بنابراین زنجیره به این صورت ساده می‌شود:
\begin{flalign*}
&\text{\lr{Encoder}}:\quad \text{\lr{Wavelet transform}} \;\to\; \text{\lr{Quantizer}} \;\to\; \text{\lr{Symbol encoder}} && \\
&\text{\lr{Decoder}}:\quad \text{\lr{Symbol decoder}} \;\to\; \text{\lr{Inverse wavelet transform}} &&
\end{flalign*}
این رویکرد، \lr{block artifacts} ناشی از افراز بلوکی اسلاید ۷۵ را برطرف می‌کند.